\documentclass[10pt,a4paper]{amsart}
\usepackage{amsmath,amssymb,amsthm,graphicx}
\usepackage{verbatim}
\newenvironment{reference}{\comment}{\endcomment}
%\newenvironment{reference}{}{}
\newenvironment{slogan}{\comment}{\endcomment}
\newenvironment{history}{\comment}{\endcomment}

% For commutative diagrams we use Xy-pic
\usepackage[all]{xy}

% We use 2cell for 2-commutative diagrams.
\xyoption{2cell}
% Theorem environments.
%
\theoremstyle{plain}
\newtheorem{theorem}[subsection]{Theorem}
\newtheorem{proposition}[subsection]{Proposition}
\newtheorem{lemma}[subsection]{Lemma}

\theoremstyle{definition}
\newtheorem{definition}[subsection]{Definition}
\newtheorem{example}[subsection]{Example}
\newtheorem{exercise}[subsection]{Exercise}
\newtheorem{situation}[subsection]{Situation}

\theoremstyle{remark}
\newtheorem{remark}[subsection]{Remark}
\newtheorem{remarks}[subsection]{Remarks}

\numberwithin{equation}{subsection}

% Macros
%
\def\lim{\mathop{\mathrm{lim}}\nolimits}
\def\colim{\mathop{\mathrm{colim}}\nolimits}
\def\Spec{\mathop{\mathrm{Spec}}}
\def\Hom{\mathop{\mathrm{Hom}}\nolimits}
\def\Ext{\mathop{\mathrm{Ext}}\nolimits}
\def\SheafHom{\mathop{\mathcal{H}\!\mathit{om}}\nolimits}
\def\SheafExt{\mathop{\mathcal{E}\!\mathit{xt}}\nolimits}
\def\Sch{\mathit{Sch}}
\def\Mor{\mathop{\mathrm{Mor}}\nolimits}
\def\Ob{\mathop{\mathrm{Ob}}\nolimits}
\def\Sh{\mathop{\mathit{Sh}}\nolimits}
\def\NL{\mathop{N\!L}\nolimits}
\def\CH{\mathop{\mathrm{CH}}\nolimits}
\def\proetale{{pro\text{-}\acute{e}tale}}
\def\etale{{\acute{e}tale}}
\def\QCoh{\mathit{QCoh}}
\def\Ker{\mathop{\mathrm{Ker}}}
\def\Im{\mathop{\mathrm{Im}}}
\def\Coker{\mathop{\mathrm{Coker}}}
\def\Coim{\mathop{\mathrm{Coim}}}

% Boxtimes
%
\DeclareMathSymbol{\boxtimes}{\mathbin}{AMSa}{"02}

%
% Macros for moduli stacks/spaces
%
\def\QCohstack{\mathcal{QC}\!\mathit{oh}}
\def\Cohstack{\mathcal{C}\!\mathit{oh}}
\def\Spacesstack{\mathcal{S}\!\mathit{paces}}
\def\Quotfunctor{\mathrm{Quot}}
\def\Hilbfunctor{\mathrm{Hilb}}
\def\Curvesstack{\mathcal{C}\!\mathit{urves}}
\def\Polarizedstack{\mathcal{P}\!\mathit{olarized}}
\def\Complexesstack{\mathcal{C}\!\mathit{omplexes}}
% \Pic is the operator that assigns to X its picard group, usage \Pic(X)
% \Picardstack_{X/B} denotes the Picard stack of X over B
% \Picardfunctor_{X/B} denotes the Picard functor of X over B
\def\Pic{\mathop{\mathrm{Pic}}\nolimits}
\def\Picardstack{\mathcal{P}\!\mathit{ic}}
\def\Picardfunctor{\mathrm{Pic}}
\def\Deformationcategory{\mathcal{D}\!\mathit{ef}}
\usepackage{cleveref}
\crefname{originallemma}{original lemma}{original lemmas}
\Crefname{originallemma}{Original lemma}{Original lemmas}
\makeatletter
\newlabel{lemma-localization-catenary}{{00NJ}{source}{}{sourcepointer-1}{}}
\newlabel{lemma-localization-catenary@cref}{{[originallemma][0][]00NJ}{[1][0][]source}{}{}{}}
\makeatother
\title{Catenarity and universal catenarity can be checked at the localisations at primes or at maximal ideals}
\author{The Stacks Project Authors}
\date{stacks-project:algebra.tex:local-snapshot-20261012; GNU FDL 1.2 or later}
\begin{document}
\maketitle
\noindent\textit{Editorial note.} All delivered mathematical text is the author's own, copied from the pinned Stacks Project chapter file algebra.tex; no statement or proof step was rewritten, reordered or removed.\par\medskip
\noindent\textit{Editorial note.} Preamble provenance, as required for this source: the frozen snapshot stages the chapter's own preamble file beside it (preamble.tex, 217 lines, 8116 bytes, the same shared Stacks-wide preamble the published chow and etale-cohomology cards use); the first line of the pinned chapter file is \textbackslash{}input\{preamble\}, so no preamble text lives inside the chapter itself. The card takes that file's lines 8-18 (the verbatim package and the three \textbackslash{}comment-based environments reference/slogan/history at lines 9-12, section 13.2, together with the Xy-pic setup) and its lines 154 to the runtime end of the file (the native theorem declarations and mathematical macros). Not one character was changed, no package was added, no macro was defined, and the excerpt serves only as preamble, never as a span coordinate.\par\medskip
\noindent\textit{Editorial note.} The same-chapter pointers are located at their \textbackslash{}label lines inside the pinned chapter file; the cross-chapter pointers are located in the project's own tag list staged in the snapshot (sources/stacks-chow/tags.txt), at the verbatim lines carrying their canonical tags.\par\medskip
\noindent\textit{Editorial note.} The counted proof refers to two same-chapter lemmas by label; both are recorded as navigation stubs pointed at their \textbackslash{}label lines, and the arguments consuming them stay inside the counted span.\par\medskip
\noindent\textit{Editorial note.} The counted spans contain no \textbackslash{}cite, so the card carries no bibliography.\par\medskip
\noindent\textit{Allowed standard prerequisites (editorial).}\begin{itemize}
\item Catenarity of localisations. Invoked in the counted proof at lines 25770-25771 and 25784-25785. Gives the implication from the ring to its localisations and supplies universal catenarity of a localisation of a localisation.
\item Correspondence of prime chains under localisation at a maximal ideal. Used in the counted proof at lines 25774-25779. Shows chains between two primes are detected after localising at a maximal ideal containing the larger one.
\end{itemize}\medskip
\ifdefined\hypertarget\hypertarget{sourcepointer-1}{}\fi
\noindent\textit{Source reference.} Original lemma 00NJ: sources/stacks-algebra/pinned-original-source.tex, source lines 25719--25719. Navigation pointer for \textbackslash{}ref\{lemma-localization-catenary\} used inside the counted span (line 25771); the argument that consumes it is retained in the span.\par\medskip
\begin{lemma}
\label{lemma-catenary-check-local}
Let $R$ be a ring. The following are equivalent
\begin{enumerate}
\item $R$ is catenary,
\item $R_\mathfrak p$ is catenary for all prime ideals $\mathfrak p$,
\item $R_\mathfrak m$ is catenary for all maximal ideals $\mathfrak m$.
\end{enumerate}
Assume $R$ is Noetherian. The following are equivalent
\begin{enumerate}
\item $R$ is universally catenary,
\item $R_\mathfrak p$ is universally catenary for all prime ideals
$\mathfrak p$,
\item $R_\mathfrak m$ is universally catenary for all maximal ideals
$\mathfrak m$.
\end{enumerate}
\end{lemma}

\begin{proof}
The implication (1) $\Rightarrow$ (2) follows from
Lemma \ref{lemma-localization-catenary} in both cases.
The implication (2) $\Rightarrow$ (3) is immediate in both cases.
Assume $R_\mathfrak m$ is catenary for all maximal ideals
$\mathfrak m$ of $R$. If $\mathfrak p \subset \mathfrak q$ are primes
in $R$, then choose a maximal ideal $\mathfrak q \subset \mathfrak m$.
Chains of primes ideals between $\mathfrak p$ and $\mathfrak q$ are
in 1-to-1 correspondence with chains of prime ideals between
$\mathfrak pR_\mathfrak m$ and $\mathfrak qR_\mathfrak m$ hence we
see $R$ is catenary. Assume $R$ is Noetherian and $R_\mathfrak m$ is
universally catenary for all maximal ideals $\mathfrak m$ of $R$.
Let $R \to S$ be a finite type ring map. Let $\mathfrak q$ be a prime
ideal of $S$ lying over the prime $\mathfrak p \subset R$.
Choose a maximal ideal $\mathfrak p \subset \mathfrak m$ in $R$.
Then $R_\mathfrak p$ is a localization of $R_\mathfrak m$ hence
universally catenary by Lemma \ref{lemma-localization-catenary}.
Then $S_\mathfrak p$ is catenary as a finite type ring over $R_\mathfrak p$.
Hence $S_\mathfrak q$ is catenary as a localization. Thus $S$ is catenary
by the first case treated above.
\end{proof}

\end{document}
